\section{Introduction}
\label{sec:intro}

Every container that runs on Kubernetes asks for resources before it starts~\citep{k8sresources}.
The numbers matter in different ways. Memory is a hard limit, so a container that
uses more than it asked for is killed and has to be run again. CPU is a soft limit, so a
container that wants more cores than it was given is throttled and finishes late. A GPU
is a choice of machine: if it is too small the task runs out of device memory, if it is
missing a GPU task fails at start, and if a CPU task gets one it pays for hardware it
never uses. In practice engineers pick these numbers by hand, usually once and usually
generously, and large clusters run well below what their jobs
reserve~\citep{verma2015borg,tirmazi2020borg}.

Autoscalers that learn from usage history~\citep{rzadca2020autopilot,k8svpa} work well for
long-running services. Batch work is harder. The first run of a batch task has no
history, and later runs often see inputs of a different size. Still, much of what
decides the right request is known before the pod starts. A scheduler can read the code
the pod will run, the size and format of its inputs, the way the pod spec stages those
inputs, and what earlier runs of the same task used. This paper asks how well different
approaches turn that information into a request, and measures the answer in dollars.

We treat the Kubernetes pod as the unit of work throughout. Teams submit pods in
different ways. The most direct is a \code{batch/v1} Job, which pairs a container image
with a manifest. Another is to write a typed Python function and let an orchestrator
build the pod. We use Flyte~\citep{flyte} as the example of the second route. Flyte acts
as an AI-native interface to Kubernetes: it turns a function into a pod, and adds typed
inputs and outputs, caching, retries and durable execution along the way. In both cases
the same user code runs in the same kind of pod. The differences are at the edges of the
pod, in how inputs arrive, what else shares the container, and what the monitoring
history reports. For Flyte, those conveniences cost 55--70\,MiB of runtime memory in the
task's container, which we measure directly (\cref{sec:validation}).

This paper makes five contributions.
\begin{enumerate}
  \item We build a benchmark of the whole pod request from 32,082 measured runs of 33
  families of generated Python tasks, covering dataframes, SQL, linear algebra,
  classical ML, PyTorch training and inference, and sharded language models. Each run
  has measured peak memory, CPU use at the 15-second resolution of cluster monitoring,
  and peak GPU memory across seven machine shapes (\cref{sec:dataset}).
  \item We present every run both as a raw Kubernetes Job and as a Flyte task. The Job
  stages its inputs in one of four ways, and its labels follow from an explicit
  accounting of what differs between the two pods (\cref{sec:twins}).
  \item We check that accounting in about 480 real pods, reading each container's own
  cgroup counters (\cref{sec:validation}). The check found two defects in our first
  release, both now fixed. The Flyte runtime constant was 2.6 times too large, and some
  init containers were too small to hold the inputs they staged. It also showed that the
  standard \code{ru\_maxrss} measurement leaves about a third of labels one cell higher
  than necessary.
  \item We score requests with a single cost model that covers memory, CPU and GPU
  together, and report both mean and dollar-weighted regret with task-clustered
  confidence intervals (\cref{sec:cost}).
  \item We compare eleven kinds of approach over three independent trainings, from static
  requests and hand-written rules to gradient-boosted trees, a zero-shot 27B model, a
  fine-tuned 1.7B decision model trained further with reinforcement learning, opinion
  pools, and a pre-registered novelty gate (\cref{sec:methods,sec:results}). Two
  controlled tests ask what the models actually read. One changes only the manifest, and
  the other removes the task's name from its code.
\end{enumerate}
